% !TEX recipe = lualatex_2
\documentclass[letterpaper,11pt]{article}

\usepackage{geometry}
\geometry{margin=1in}

\usepackage{fontspec}
\setmainfont[Ligatures=TeX]{Latin Modern Roman}

\usepackage{hyperref}
\hypersetup{colorlinks=true, urlcolor=black}

\usepackage{parskip}
\setlength{\parskip}{0.8em}
\setlength{\parindent}{0pt}

\pagestyle{empty}

\begin{document}

\textbf{Rodrigo G. Freundt}\\
211 S Titus Ave, Ithaca, NY 14850\\
rodrigo.freundt@gmail.com\\

\vspace{1em}
\today

\vspace{1em}
Hiring Committee\\
Smithsonian Astrophysical Observatory (SAO)\\
Cambridge, MA

\vspace{1em}
Dear Hiring Committee,

I am writing to apply for the Electronics Engineer position (applying concurrently to both IS-0855-14 and IS-0855-13) at the Smithsonian Astrophysical Observatory. With over ten years of experience designing, deploying and maintaining state-of-the-art astronomical instrumentation, I believe I am exceptionally well suited for this role.

I am currently a Ph.D. candidate in Astronomy and Space Sciences at Cornell University. I work with Prof. Gordon Stacey and Prof. Michael Niemack on multiple aspects of Prime-Cam, a cryogenic modular receiver for the Fred Young Submillimeter Telescope (FYST). In particular, I've been focused on the development of the Epoch of Reionization Spectrometer (EoR-Spec), one of Prime-Cam's instrument modules.

I got my professional degree in Electronics Engineering from the Pontifical Catholic University of Peru, where my research focused on radio astronomy instrumentation. Before joining Cornell, I worked for two years as site engineer at the Atacama Cosmology Telescope (ACT), a cosmology experiment that was located in northern Chile.

%I am currently a Ph.D. candidate in Astronomy and Space Sciences at Cornell University, where I lead the electronics and systems development for EoR-Spec, a direct-detection imaging spectrometer for the CCAT Observatory's Prime-Cam receiver. In this role I have been responsible for the full engineering lifecycle: from translating science requirements into hardware concepts, through schematic and PCB design, fabrication oversight, cryogenic integration, and on-site commissioning. I have designed detector readout and thermometry harnesses operating at millikelvin temperatures, led design reviews within a large multi-institutional collaboration, and authored technical reports and peer-reviewed conference proceedings documenting this work. This experience maps directly onto the Major Duties described in your posting.

%Prior to my doctoral work, I served as Site Engineer for the Atacama Cosmology Telescope (ACT) in the Chilean Atacama Desert, where I was responsible for the continuous operation of a complex cryogenic receiver at 5,200 meters altitude. This included maintenance and troubleshooting of electrical, cryogenic, and mechanical systems; high-precision photogrammetry alignment; Fourier Transform Spectroscopy passband characterization; and support for detector array upgrades. Operating in a resource-constrained, remote environment sharpened my ability to diagnose and resolve complex technical problems independently and under pressure — a skill I expect to be equally valuable at the CfA.

%My technical toolkit aligns well with your requirements. I have hands-on experience with KiCAD for schematic capture and PCB layout, FPGA-based digital signal processing systems (developed during my undergraduate thesis on a polyphase filter bank spectrometer using FlexRIO FPGA and CUDA-enabled GPU), and a working knowledge of thermal and mechanical engineering constraints as they apply to instrument packaging. I am also proficient in Python, C/C++, and LabVIEW, and regularly use simulation and modeling tools in my design process.

The prospect of contributing to the SAO is genuinely exciting to me. I would welcome the opportunity to discuss how my background can support SAO's projects and goals.

Thank you for your consideration.

\vspace{1.5em}
Sincerely,

\vspace{2em}
Rodrigo G. Freundt

\vspace{2em}
Notes: CV attached as a separate file; Expected completion of Ph.D. degree: Spring 2027; Authorized to work in the US for any employer.

\end{document}